\documentclass[11pt,letterpaper]{article}

% CODIFICACIÓN E IDIOMA
\usepackage[utf8]{inputenc}     
\usepackage[T1]{fontenc}         % Salida de fuentes en T1: acentos y guionado correctos
\usepackage[spanish]{babel}      % Localización: "Tabla", "Figura", "Demostración", fechas, etc.
\decimalpoint
\usepackage{microtype}           % Micro-ajustes tipográficos (mejora notable de legibilidad)

% MATEMÁTICAS

\usepackage{amsmath,amssymb,amsfonts,amsthm}
\usepackage{mathtools}
\providecommand{\abs}[1]{\lvert#1\rvert}
\providecommand{\norm}[1]{\lVert#1\rVert}
\newcommand{\Aut}{\operatorname{Aut}}

% TIPOGRAFÍA

\usepackage{newpxtext,newpxmath}
\usepackage{titlesec}            % Formato de \section, \subsection, etc.

% GEOMETRÍA DE PÁGINA
\usepackage[a4paper,letterpaper,margin=2.5cm,headheight=14.5pt,includeheadfoot]{geometry}



% FIGURAS, TABLAS Y CÓDIGO

\usepackage{graphicx}
\graphicspath{{imagenes/}{figuras/}{img/}}
\usepackage{subcaption}          % Sustituye a "subfigure" (paquete obsoleto desde 2001)
\usepackage{caption}
\captionsetup{font=small,labelfont=bf,margin=8pt}
\usepackage{booktabs}            % Tablas profesionales (\toprule, \midrule, \bottomrule)
\usepackage{array}
\usepackage{siunitx}
% \sisetup{output-decimal-marker = {.}, group-separator = {,}}
% \sisetup{
%     locale = US,
%     output-decimal-marker = {.},
%     group-separator = {,}
% }
\usepackage{enumitem}            % Listas con espaciado controlado
\usepackage{listings}
\usepackage{xcolor}
\usepackage{tikz}
\usepackage{pgfplots}
\pgfplotsset{compat=1.18}        % Versión de compatibilidad actualizada (1.15 → 1.18)
\usetikzlibrary{arrows.meta}     % Sustituye a la librería "arrows", ya obsoleta


% ESPACIADO DE PÁRRAFOS
\setlength{\parindent}{0pt}
\setlength{\parskip}{0.65em plus 0.1em minus 0.1em}


% HIPERVÍNCULOS (se carga casi al final, como recomienda su documentación)
\usepackage{hyperref}
\hypersetup{
    colorlinks   = true,
    linkcolor    = {primary},
    citecolor    = {primary},
    urlcolor     = {primary!70!black},
    pdfborder    = {0 0 0},
    breaklinks   = true,
    bookmarksnumbered = true
}


% PALETA DE COLORES

\definecolor{primary}{RGB}{120,25,90}    % Vino/púrpura institucional (evolución de "purpura")
\definecolor{secondary}{RGB}{70,70,70}   % Gris oscuro para texto secundario
\definecolor{bgcode}{RGB}{247,247,250}   % Fondo suave para bloques de código
\hypersetup{linkcolor=primary, citecolor=primary, urlcolor=primary!70!black}

% TÍTULOS DE SECCIÓN 

\titleformat{\section}
  {\normalfont\Large\bfseries\sffamily\color{primary}}
  {\thesection}{0.8em}{}[{\color{primary!60}\titlerule[0.8pt]}]
\titlespacing*{\section}{0pt}{2.8ex plus 1ex minus .2ex}{1.4ex plus .2ex}

\titleformat{\subsection}
  {\normalfont\large\bfseries\sffamily\color{primary!85!black}}
  {\thesubsection}{0.8em}{}
\titlespacing*{\subsection}{0pt}{2.2ex plus 1ex minus .2ex}{1ex plus .2ex}

\titleformat{\subsubsection}
  {\normalfont\normalsize\bfseries\sffamily\color{secondary}}
  {\thesubsubsection}{0.8em}{}


% ENCABEZADO Y PIE DE PÁGINA

\usepackage{fancyhdr}
\pagestyle{fancy}
\fancyhf{}
\renewcommand{\headrulewidth}{0.6pt}
\renewcommand{\footrulewidth}{0.4pt}
\fancyheadoffset{0pt}

% Estilo "plain" (usado en la portada) sin encabezado
\fancypagestyle{plain}{%
  \fancyhf{}
  \renewcommand{\headrulewidth}{0pt}
  \renewcommand{\footrulewidth}{0pt}
}

% Datos del documento 
\newcommand{\asignaturaNombre}{Matemáticas}
\newcommand{\tareaNombre}{Tarea 3}
\newcommand{\autorNombre}{Eduardo Alanís García}
\newcommand{\institucionNombre}{MCDI}
\newcommand{\profesorNombre}{Dra. Briceyda B. Delgado López}
\newcommand{\fechaEntrega}{29 de octubre de 2026}


\fancyhead[L]{\small\color{primary}\textbf{\asignaturaNombre}}
\fancyhead[R]{\small\color{primary}\textbf{\tareaNombre}}
\fancyfoot[L]{\small\color{secondary}\autorNombre}
\fancyfoot[C]{\small\color{secondary}\thepage}
\fancyfoot[R]{\small\color{secondary}\institucionNombre}


% ENTORNOS TIPO TEOREMA
\newtheoremstyle{coloredplain}
  {0.9em}{0.9em}{\normalfont}{0pt}{\bfseries\color{primary}}{.}{0.6em}{}
\newtheoremstyle{coloreddef}
  {0.9em}{0.9em}{\normalfont}{0pt}{\bfseries\color{primary!75!black}}{.}{0.6em}{}

\theoremstyle{coloredplain}
\newtheorem{teorema}{Teorema}[section]
\newtheorem{lema}[teorema]{Lema}
\newtheorem{proposicion}[teorema]{Proposición}
\newtheorem{corolario}[teorema]{Corolario}

\theoremstyle{coloreddef}
\newtheorem{definicion}{Definición}[section]
\newtheorem{ejemplo}[definicion]{Ejemplo}
\newtheorem{observacion}[definicion]{Nota}
\newtheorem{afirmacion}[definicion]{Afirmación}
\newtheorem{pregunta}[definicion]{Pregunta}
\newtheorem{respuesta}[definicion]{Respuesta}
\newtheorem{ejercicio}{Ejercicio}


\renewcommand{\qedsymbol}{$\blacksquare$}
\newenvironment{solucion}
  {\par\noindent\textbf{\color{primary}Solución.}\ }
  %{\hfill$\blacksquare$\par}

\newenvironment{demostraccion}
  {\par\noindent\textbf{\color{primary}Demostración.}\ }
  {\hfill$\blacksquare$\par}

\newcommand{\pie}[1]{\footnote{#1}}


% CÓDIGO FUENTE (bloques "listings" con estilo unificado)

\lstdefinestyle{Rstyle}{
    language        = R,
    backgroundcolor = \color{bgcode},
    basicstyle      = \ttfamily\small,
    keywordstyle    = \color{primary}\bfseries,
    commentstyle    = \color{secondary}\itshape,
    stringstyle     = \color{teal},
    numbers         = left,
    numberstyle     = \tiny\color{secondary},
    numbersep       = 8pt,
    frame           = single,
    rulecolor       = \color{primary!30},
    breaklines      = true,
    showstringspaces= false,
    captionpos      = b,
    tabsize         = 2
}
\lstset{style=Rstyle}
\lstset{literate={á}{{\'a}}1 {é}{{\'e}}1 {í}{{\'i}}1 {ó}{{\'o}}1 {ú}{{\'u}}1 {ñ}{{\~n}}1 {Á}{{\'A}}1 {É}{{\'E}}1 {Í}{{\'I}}1 {Ó}{{\'O}}1 {Ú}{{\'U}}1}

% Referencias
\usepackage{csquotes}
\usepackage[numbers]{natbib}
% ============================================================================


% ============================================================================
\begin{document}


% PORTADA

\begin{titlepage}
  \thispagestyle{plain}
  \centering
  \vspace*{2.5cm}

  {\sffamily\color{secondary}\large Maestría en Ciencia de Datos e Información \par}
  \vspace{0.3cm}
  {\color{primary}\rule{0.5\textwidth}{0.8pt}\par}
  \vspace{2.2cm}

  {\sffamily\bfseries\color{primary}\Huge \tareaNombre \par}
  \vspace{0.6cm}
  {\sffamily\Large \asignaturaNombre \par}

  \vspace{3.5cm}
  {\color{primary!60}\rule{0.6\textwidth}{0.4pt}\par}
  \vspace{1cm}

  \begin{tabular}{rl}
    \textbf{Alumno:}   & \autorNombre \\[2pt]
    \textbf{Profesora:} & \profesorNombre \\[2pt]
    \textbf{Entrega:}    & \fechaEntrega \\
  \end{tabular}

  \vfill
  {\small\color{secondary} INFOTEC \par}
\end{titlepage}

% Descomente si desea un índice automático:
% \tableofcontents
% \newpage


\section*{Instrucciones}
Los procedimientos se desarrollan de forma analítica y ordenada. El código se utiliza únicamente en el ejercicio 6, donde está permitido.

\section*{Herramientas utilizadas}
Todos los ejercicios se desarrollaron de manera analítica y el documento se elaboró con \LaTeX{}. El único apartado en el que se permitió el uso de código, el Ejercicio 6, se realizó con \textbf{R}.


\section*{Ejercicios}

\begin{ejercicio}[15 puntos] Considere la función $f(x,y)=x^4+y^4-4xy+2x^2+2y^2$.

\begin{enumerate}[label=(\alph*)]
  \item Calcule el gradiente de $f$.

  \begin{solucion}
  Derivando parcialmente con respecto a $x$ y a $y$ se obtiene
  \[
    \nabla f(x,y)=
    \begin{pmatrix}
      4x^3-4y+4x\\
      4y^3-4x+4y
    \end{pmatrix}.
  \] \citetext{\citealp[diapositiva~5]{delgado2026}; \citealp[Secciones~14.3 y 14.6]{stewart2012}}.
  \end{solucion}




  
  \item Calcule la matriz Hessiana de $f$.

  \begin{solucion}
      Las segundas derivadas parciales son
  \[
    f_{xx}=12x^2+4,\qquad f_{xy}=f_{yx}=-4,\qquad f_{yy}=12y^2+4,
  \]
  por lo tanto, la matriz Hessiana es
  \[
    H_f(x,y)=
    \begin{pmatrix}
      12x^2+4 & -4\\
      -4 & 12y^2+4
    \end{pmatrix}.
  \] \citetext{\citealp[diapositiva~6]{delgado2026}; \citealp[Sección~14.7]{stewart2012}}.
  \end{solucion}




  
  \item Determine los puntos críticos de $f$.

  \begin{solucion}
      Para encontrar los puntos críticos igualamos el gradiente a cero \citep[diapositiva~5]{delgado2026}:
  \[
    \begin{cases}
      4x^3-4y+4x=0,\\
      4y^3-4x+4y=0.
    \end{cases}
    \quad\Longleftrightarrow\quad
    \begin{cases}
      y=x^3+x,\\
      x=y^3+y.
    \end{cases}
  \]
  Al sumar las ecuaciones $x^3=y-x$ y $y^3=x-y$, obtenemos
  \[
    x^3+y^3=0
    \quad\Longleftrightarrow\quad
    (x+y)(x^2-xy+y^2)=0.
  \]
  Si $x^2-xy+y^2=0$, entonces $x=y=0$. En el otro caso, $y=-x$ y, usando $y=x^3+x$,
  \[
    -x=x^3+x
    \quad\Longleftrightarrow\quad
    x(x^2+2)=0.
  \]
  Como $x^2+2>0$ $\forall x \in \mathbb{R}$, se obtiene $x=0$ y, por consiguiente, $y=0$. El único punto crítico es $(0,0)$.
  \end{solucion}

  
  \item Analice cada punto crítico utilizando la matriz Hessiana y clasifíquelo como mínimo, máximo o punto de silla.

  \begin{solucion}
       El único punto crítico es $(0,0)$, en ese punto la Hessiana vale
  \[
    H_f(0,0)=
    \begin{pmatrix}
      4 & -4\\
      -4 & 4
    \end{pmatrix},
    \qquad \det H_f(0,0)=0.
  \]
  Sus valores propios son $0$ y $8$, por lo que la Hessiana es semidefinida positiva, pero el criterio de la Hessiana no permite clasificar el punto porque su determinante es cero. Para concluir, reescribimos la función:
  \[
    f(x,y)= x^4+y^4-4xy+2x^2+2y^2 = x^4+y^4+2(x-y)^2\geq 0.
  \]
  La igualdad se cumple únicamente cuando $x=y=0$. Por lo tanto, $(0,0)$ es un mínimo estricto global \citetext{\citealp[diapositivas~7--10]{delgado2026}; \citealp[Sección~14.7]{stewart2012}}.
  \end{solucion}
  
\end{enumerate}
\end{ejercicio}








\begin{ejercicio}[20 puntos] Considere la función 
\[
  f(x,y)=x^2+xy+y^2+3x
\]
\begin{enumerate}[label=(\alph*)]
  \item Encuentra el mínimo (o mínimos) analíticamente. ¿Son mínimos globales o relativos?

  \begin{solucion}
   El gradiente y la Hessiana son
    \[
      \nabla f(x,y)=
      \begin{pmatrix}2x+y+3\\x+2y\end{pmatrix},
      \qquad
      H_f(x,y)=
      \begin{pmatrix}2&1\\1&2\end{pmatrix}.
    \]
    Al resolver $\nabla f=\vec{0}$, de $x+2y=0$ obtenemos $x=-2y$. Sustituyendo en $2x+y+3=0$,
    \[
      -4y+y+3=0
      \quad\Longrightarrow\quad y=1,\qquad x=-2.
    \]
    Los valores propios de la Hessiana son $1$ y $3$, ambos positivos \citep[diapositivas~7--10]{delgado2026}. Como la Hessiana es
constante, es definida positiva en todo $\mathbb{R}^2$ y, por lo tanto, $f$ es estrictamente
convexa \citep[Capítulo~2]{nocedal2006}. Así, el único punto crítico $(-2,1)$ es un mínimo global, y su valor es
\[
  f(-2,1)=4-2+1-6=-3.
\]


Otra forma de verificar que se trata de un mínimo global es expresar la función de la siguiente manera:
\[
\begin{aligned}
f(x,y) &= x^2 + xy + y^2 + 3x \\[6pt]
       &= x^2 + x(y+3) + y^2 \\[6pt]
       &= \left(x+\frac{y+3}{2}\right)^2
          -\frac{(y+3)^2}{4} + y^2 \\[6pt]
       &= \left(x+\frac{y+3}{2}\right)^2
          +\frac{4y^2-(y+3)^2}{4} \\[6pt]
       &= \left(x+\frac{y+3}{2}\right)^2
          +\frac{3}{4}(y^2-2y-3) \\[6pt]
       &= \left(x+\frac{y+3}{2}\right)^2
          +\frac{3}{4}(y^2-2y+1-4) \\[6pt]
       &= \left(x+\frac{y+3}{2}\right)^2
          +\frac{3}{4}\left[(y-1)^2-4\right] \\[6pt]
       &= \left(x+\frac{y+3}{2}\right)^2
          +\frac{3}{4}(y-1)^2-3 \\[6pt]
       &= -3+\left(x+\frac{y+3}{2}\right)^2
          +\frac{3}{4}(y-1)^2.
\end{aligned}
\]

Como los términos cuadrados son no negativos, se cumple que
\[
f(x,y)\geq -3
\qquad \text{para todo } (x,y)\in\mathbb{R}^2.
\]

Por lo tanto, la función tiene un mínimo global único en el punto \((-2,1)\), cuyo valor mínimo es \(f(-2,1)=-3\).

  \end{solucion}
  
  \item Construye cuatro curvas de nivel de $f(x,y)$ [líneas de valor constante de $f(x,y)$].

  \begin{solucion}
      Utilizando la expresión obtenida en el inciso anterior, podemos escribir
\begin{align*}
f(x,y)
&= -3+\left(x+2+\frac{y-1}{2}\right)^2
+\frac{3}{4}(y-1)^2\\
&= -3+(x+2)^2+(x+2)(y-1)
+\frac{(y-1)^2}{4}+\frac{3}{4}(y-1)^2\\
&= -3+(x+2)^2+(x+2)(y-1)+(y-1)^2.
\end{align*}

Por lo tanto, la función puede expresarse de la siguiente manera:

\[
f(x,y)=x^2+xy+y^2+3x
=-3+(x+2)^2+(x+2)(y-1)+(y-1)^2.
\]

Las curvas de nivel de la función están dadas por la ecuación \citep[Sección~14.1]{stewart2012}
\(f(x,y)=c\). Al sustituir la expresión anterior, obtenemos
\[
(x+2)^2+(x+2)(y-1)+(y-1)^2=c+3.
\]

Como el valor mínimo de la función es \(f(-2,1)=-3\), para obtener
curvas de nivel debemos elegir valores de \(c>-3\).
En particular, si tomamos \(c=-2,-1,0,4\), obtenemos las siguientes
ecuaciones:
    \begin{align*}
      (x+2)^2+(x+2)(y-1)+(y-1)^2&=1,\\
      (x+2)^2+(x+2)(y-1)+(y-1)^2&=2,\\
      (x+2)^2+(x+2)(y-1)+(y-1)^2&=3,\\
      (x+2)^2+(x+2)(y-1)+(y-1)^2&=7.
    \end{align*}
    
  \end{solucion}


  
  \item ¿Es la búsqueda univariante un buen método numérico para encontrar el óptimo de $f(x,y)$? Justifique su respuesta.

  \begin{solucion}
      Sí, es un buen método, porque $f$ es cuadrática y está bien escalada: su Hessiana $H_f=\begin{pmatrix}2&1\\1&2\end{pmatrix}$ tiene valores propios $1$ y $3$, cuyo cociente es $3$, un valor pequeño, por lo que las curvas de nivel son elipses poco alargadas. Lo importante es elegir bien las direcciones de búsqueda \citep[Secciones~5.2--5.3]{tierney2023}.
  \end{solucion}
  \item Supongamos que la dirección de búsqueda viene dada por $\vec{s}=[1,0]$. Si inicializamos en el punto $(0,0)$, encuentra el punto óptimo $P_1$ en esa dirección de búsqueda analíticamente, no numéricamente.

\begin{solucion}
     La recta de búsqueda que pasa por $(0,0)$ con dirección $\vec{s}=[1,0]$ es
$(x,y)=(\lambda,0)$ con $\lambda\in\mathbb{R}$, pues el óptimo se busca sobre toda la recta
    \citep[Sección~5.2]{tierney2023}:
    \[
      f(\lambda,0)=\lambda^2+3\lambda.
    \]
    Derivando e igualando a cero,
    \[
      2\lambda+3=0
      \quad\Longrightarrow\quad
      \lambda=-\frac{3}{2}.
    \]
    Por lo tanto, $P_1=(-\frac{3}{2},0)$.
\end{solucion}
  
  \item Repita el ejercicio para un punto de partida de $(0,4)$ para encontrar $P_2$.

  \begin{solucion}
    La recta de búsqueda es $(x,y)=(\lambda,4)$. En ella,
    \[
      f(\lambda,4)=\lambda^2+7\lambda+16.
    \]
    De $2\lambda+7=0$ resulta $\lambda=-\frac{7}{2}$; por tanto, $P_2=(-\frac{7}{2},4)$.
  \end{solucion}


  
  \item Demuestre gráficamente que una recta que une $P_1$ y $P_2$, pasa por el óptimo.


  \begin{solucion}
    La pendiente de la recta que pasa por $P_1=(-\frac{3}{2},0)$ y $P_2=(-\frac{7}{2},4)$ es
    \[
      m=\frac{4-0}{-\frac{7}{2}+\frac{3}{2}}=-2,
    \]
    así que su ecuación es $y=-2x-3$. El óptimo $(-2,1)$ satisface esta ecuación. La figura muestra las cuatro curvas de nivel, los puntos $P_1$ y $P_2$, la recta que los une y el mínimo.

        \begin{center}
      \begin{tikzpicture}[scale=0.82]
        % Curvas de nivel: elipses centradas en (-2,1), con eje mayor a -45 grados.
        \draw[blue!55,rotate around={-45:(-2,1)}] (-2,1) ellipse [x radius=3.742,y radius=2.160];
        \draw[blue!55,rotate around={-45:(-2,1)}] (-2,1) ellipse [x radius=2.449,y radius=1.414];
        \draw[blue!55,rotate around={-45:(-2,1)}] (-2,1) ellipse [x radius=2.000,y radius=1.155];
        \draw[blue!55,rotate around={-45:(-2,1)}] (-2,1) ellipse [x radius=1.414,y radius=0.816];
        \draw[->,gray] (-5.5,0)--(1.35,0) node[right] {$x$};
        \draw[->,gray] (0,-2.35)--(0,4.35) node[above] {$y$};
        \draw[red!75!black,dashed,thick] (-1.5,0)--(-3.5,4);
        \fill[black] (-1.5,0) circle (2pt) node[below right] {$P_1$};
        \fill[black] (-3.5,4) circle (2pt) node[above left] {$P_2$};
        \fill[red!75!black] (-2,1) circle (2.5pt) node[below right] {$(-2,1)$};
      \end{tikzpicture}
    \end{center}
  \end{solucion}
  
\end{enumerate}
    
\end{ejercicio}




\begin{ejercicio}[15 puntos] Dada la función $h(x,y,z)=x^2+y^2+z^2$ sujeta a la restricción $x+y+z=1$:
\begin{enumerate}[label=(\alph*)]
  \item Encuentre los puntos estacionarios utilizando el método de multiplicadores de Lagrange.

\begin{solucion}
    Definimos el Lagrangiano
    \[
      L(x,y,z,\lambda)=x^2+y^2+z^2-\lambda(x+y+z-1).
    \]
    Las derivadas parciales son
    \[
      \begin{cases}
        \frac{\partial L}{\partial x} = 2x-\lambda=0,\\
        \frac{\partial L}{\partial y} =2y-\lambda=0,\\
       \frac{\partial L}{\partial z} = 2z-\lambda=0,\\
        \frac{\partial L}{\partial \lambda} = -(x+y+z-1) = 0.
      \end{cases}
    \]
    De las primeras tres ecuaciones se obtiene $x=y=z=\lambda/2$. Sustituyendo en la restricción,
    \[
      3\frac{\lambda}{2}=1
      \quad\Longrightarrow\quad \lambda=\frac23,
      \qquad x=y=z=\frac13.
    \]
    El único punto estacionario es $(\frac13,\frac13,\frac13)$, con $\lambda = \frac{2}{3}$ \citetext{\citealp[diapositivas~11--13]{delgado2026}; \citealp[Sección~14.8]{stewart2012}}.
\end{solucion}

  
  \item Verifique si los puntos encontrados son máximos o mínimos utilizando la Hessiana del Lagrangiano.

  \begin{solucion}
La Hessiana del Lagrangiano respecto de $(x,y,z)$ es
    \[
      H_L=
      \begin{pmatrix}
        2&0&0\\
        0&2&0\\
        0&0&2
      \end{pmatrix},
    \]
    y el gradiente de la restricción es $\nabla g=(1,1,1)$. La matriz Hessiana orlada es
    \[
      B=
      \begin{pmatrix}
        0 & \nabla g \\
        \nabla g^{\top} & H_L
      \end{pmatrix}
      =
      \begin{pmatrix}
        0 & 1 & 1 & 1 \\
        1 & 2 & 0 & 0 \\
        1 & 0 & 2 & 0 \\
        1 & 0 & 0 & 2
      \end{pmatrix}.
    \]

   Sea $\Delta_2^B$ el determinante del bloque formado por las tres primeras filas y columnas, y $\Delta_3^B=\det B$:

    \[
\Delta_2^B =
\begin{vmatrix}
0 & 1 & 1 \\
1 & 2 & 0 \\
1 & 0 & 2
\end{vmatrix}
= -4 < 0,
\qquad
\Delta_3^B =
\begin{vmatrix}
0 & 1 & 1 & 1 \\
1 & 2 & 0 & 0 \\
1 & 0 & 2 & 0 \\
1 & 0 & 0 & 2
\end{vmatrix}
= -12 < 0.
\]

El criterio establece que si $\Delta_k^B<0$ para todo $k$, es decir, si todos los menores orlados son negativos, el punto es un \emph{mínimo local restringido} \citep[diapositivas~14--16]{delgado2026}. Como $\Delta_2^B<0$ y $\Delta_3^B<0$, concluimos que
    \[
      P=\left(\tfrac13,\tfrac13,\tfrac13\right)
    \]
    es un mínimo local restringido, con valor $h(P)=3\cdot\tfrac19=\tfrac13$.   Además, como $f$ es convexa y la restricción es lineal, este minimo es \emph{global} \citep[diapositiva~22]{delgado2026}.
    
  \end{solucion}
\end{enumerate}
    
\end{ejercicio}



\begin{ejercicio}[10 puntos] Aplique el método de gradiente descendente para minimizar la función
\[
  f(x,y)=x^2+2y^2-xy-6x-8y
\]
partiendo del punto inicial $(x_0,y_0)=(0,0)$.
\begin{enumerate}[label=(\alph*)]
  \item Use un tamaño de paso constante $\alpha=0.1$.
\begin{solucion}
El gradiente de la función es
\[
  \nabla f(x,y)=
  \begin{pmatrix}
    2x-y-6\\
    4y-x-8
  \end{pmatrix}.
\]
La actualización del gradiente descendente es
\[
  \begin{aligned}
    x_{k+1}&=x_k-0.1(2x_k-y_k-6)=0.8x_k+0.1y_k+0.6,\\
    y_{k+1}&=y_k-0.1(4y_k-x_k-8)=0.1x_k+0.6y_k+0.8.
  \end{aligned}
\]
\end{solucion}

  \item Realice al menos 5 iteraciones y analice la convergencia.

  \begin{solucion}
      Aplicando estas expresiones desde $(0,0)$, se obtienen las siguientes iteraciones:
\[
\begin{array}{c|r|r|r}
  k & x_k & y_k & f(x_k,y_k)\\
  \hline
  0 & 0.00000 & 0.00000 & 0.000000\\
  1 & 0.60000 & 0.80000 & -8.840000\\
  2 & 1.16000 & 1.34000 & -14.297600\\
  3 & 1.66200 & 1.72000 & -17.911596\\
  4 & 2.10160 & 1.99820 & -20.392288\\
  5 & 2.48110 & 2.20908 & -22.124262
\end{array}
\]

Para identificar el punto al que converge el método, igualamos el gradiente a cero:
\[
  \begin{cases}
    2x-y-6=0,\\
    4y-x-8=0.
  \end{cases}
  \quad\Longrightarrow\quad
  (x^*,y^*)=\left(\frac{32}{7},\frac{22}{7}\right),
  \qquad f(x^*,y^*)=-\frac{184}{7} \approx -26.2857
\]
La Hessiana constante es
\[
  H_f=\begin{pmatrix}2&-1\\-1&4\end{pmatrix},
\]
con valores propios $3-\sqrt{2}$ y $3+\sqrt{2}$. Como ambos son positivos, la Hessiana es
definida positiva y $f$ es estrictamente convexa; en consecuencia,
$z^*=(x^*,y^*)=(32/7,22/7)$ es el mínimo global.

Para una función cuadrática estrictamente convexa, el gradiente descendente con paso
constante converge si $0<\alpha<2/\lambda_{\max}$ \citep[Sección~5.13]{tierney2023}. Aquí
\[
  0.1<\frac{2}{3+\sqrt2}\approx0.453,
\]
así que las iteraciones convergen al mínimo global. Esto es consistente con la tabla, donde
$f(x_k,y_k)$ disminuye en cada iteración.

  \end{solucion}
  
\end{enumerate}
    
\end{ejercicio}










\begin{ejercicio}[15 puntos] Verifique que la matriz Hessiana de la siguiente función
\[
  f(x,y)=(1.5-x(1-y))^2+(2.25-x(1-y^2))^2+(2.625-x(1-y^3))^2,
\]
no es positiva definida en una vecindad del punto $(0,1)$.

\begin{solucion}
Para abreviar las operaciones, consideramos
\[
\begin{aligned}
r_1(x,y)&=1.5-x(1-y),\\
r_2(x,y)&=2.25-x(1-y^2),\\
r_3(x,y)&=2.625-x(1-y^3).
\end{aligned}
\]
Así,
\[
f(x,y)=r_1^2+r_2^2+r_3^2=\sum_{i=1}^{3} r_i^2 .
\]

Derivando $r_1,r_2,r_3$ respecto a $x$ y a $y$:
\[
\begin{aligned}
\frac{\partial r_1}{\partial x}&=y-1,
&\frac{\partial r_1}{\partial y}&=x,\\
\frac{\partial r_2}{\partial x}&=y^2-1,
&\frac{\partial r_2}{\partial y}&=2xy,\\
\frac{\partial r_3}{\partial x}&=y^3-1,
&\frac{\partial r_3}{\partial y}&=3xy^2.
\end{aligned}
\]

Por la regla de la cadena,
$\frac{\partial}{\partial x}(r_i^2)=2r_i\,\frac{\partial r_i}{\partial x}$, y lo
mismo para $y$. Entonces
\[
f_x(x,y)=2r_1(y-1)+2r_2(y^2-1)+2r_3(y^3-1).
\]
Para $f_y$:
\[
f_y(x,y)=2r_1\,x+2r_2\,(2xy)+2r_3\,(3xy^2)
=2x\left(r_1+2y\,r_2+3y^2\,r_3\right).
\]

Sustituyendo $x=0$, $y=1$:
\[
\begin{aligned}
r_1(0,1)&=1.5,\\
r_2(0,1)&=2.25,\\
r_3(0,1)&=2.625.
\end{aligned}
\]
Además, $y-1=y^2-1=y^3-1=0$ y $x=2xy=3xy^2=0$; es decir, todas las derivadas
parciales de $r_1,r_2,r_3$ se anulan en $(0,1)$. En consecuencia,
$f_x(0,1)=0$ y $f_y(0,1)=0$, o sea,
\[
\nabla f(0,1)=
\begin{pmatrix}
0\\
0
\end{pmatrix}.
\]
Por lo tanto, $(0,1)$ es un punto estacionario.


\emph{Derivada $f_{xx}$.} Los factores $(y-1)$, $(y^2-1)$ y $(y^3-1)$ no
dependen de $x$, y $\frac{\partial r_i}{\partial x}$ es justamente cada uno de
ellos:
\begin{align*}
f_{xx}(x,y)
&= 2(y-1)(y-1) + 2(y^2-1)(y^2-1) + 2(y^3-1)(y^3-1)\\
&= 2\left[(y-1)^2 + (y^2-1)^2 + (y^3-1)^2\right].
\end{align*}

\emph{Derivada $f_{yy}$.} Sea $g=r_1+2y\,r_2+3y^2\,r_3$, de modo que
$f_y=2x\,g$. Como $2x$ no depende de $y$, se tiene $f_{yy}=2x\,g_y$. Al derivar
$g$ respecto a $y$ hay que usar la regla del producto en los términos
$2y\,r_2$ y $3y^2\,r_3$:
\[
g_y=x+2r_2+2y(2xy)+6y\,r_3+3y^2(3xy^2).
\]
Por lo tanto,
\[
f_{yy}(x,y)=2x\left[x+2r_2+2y(2xy)+6y\,r_3+3y^2(3xy^2)\right].
\]

\emph{Derivada $f_{xy}$.} Derivamos $f_x$ respecto a $y$ con la regla del
producto:
\begin{align*}
f_{xy}(x,y)
&= 2\bigl[x(y-1)+r_1\bigr]+2\bigl[2xy(y^2-1)+2y\,r_2\bigr]
   +2\bigl[3xy^2(y^3-1)+3y^2\,r_3\bigr]\\
&= 2\left(r_1 + 2y\,r_2 + 3y^2\,r_3\right)
+ 2x\left[(y-1) + 2y(y^2-1) + 3y^2(y^3-1)\right].
\end{align*}

Evaluar en $(0,1)$.
\begin{align*}
f_{xx}(0,1) &= 2[0+0+0]=0,\\
f_{yy}(0,1) &= 0 \qquad (\text{pues el factor } 2x \text{ vale } 0),\\
f_{xy}(0,1) &= 2\left(1\cdot 1.5 + 2\cdot 2.25 + 3\cdot 2.625\right)
             = 2(1.5+4.5+7.875)=2(13.875)=\frac{111}{4}.
\end{align*}
En $f_{xy}$, el segundo sumando se anula porque $x=0$. Con estos valores,
\[
H_f(0,1)=
\begin{pmatrix}
0&\frac{111}{4}\\[2pt]
\frac{111}{4}&0
\end{pmatrix},
\qquad
\det H_f(0,1)=0\cdot 0-\left(\frac{111}{4}\right)^2
=-\left(\frac{111}{4}\right)^2<0.
\]
Los valores propios cumplen $\lambda^2-\left(\frac{111}{4}\right)^2=0$, es
decir, $\lambda=\pm\frac{111}{4}$. Hay uno positivo y otro negativo, así que $H_f(0,1)$ es indefinida y el punto $(0,1)$ es un punto silla \citep[diapositivas~7--10]{delgado2026}.

La función $\det H_f(x,y)=f_{xx}f_{yy}-f_{xy}^2$ es un polinomio en $x,y$, por lo que es continua. Como es negativa en $(0,1)$, la continuidad garantiza que conserva ese signo cerca de $(0,1)$. % Existe, entonces, $\delta>0$ tal que
% $\det H_f(x,y)<0$ para todo $(x,y)$ a distancia menor que $\delta$ de $(0,1)$.
Es decir, existe una vecindad de $(0,1)$ en la que el determinante permanece negativo. Lo que implica, la Hessiana no es positiva definida en esa vecindad.
\end{solucion}



\begin{enumerate}[label=(\alph*)]
  \item ¿Qué pasará con el método de Newton si lo tomamos como punto inicial?
  
  \begin{solucion}
El método de Newton para optimización consiste en aplicar la iteración de Newton
para sistemas de ecuaciones $F(x)=0$ a $F(x)=\nabla f(x)$, tomando como
Jacobiano a la Hessiana $H_f$.  Newton genera una sucesión de puntos mediante
  \[
  \mathbf{x}_{k+1}
  =\mathbf{x}_k-H_f(\mathbf{x}_k)^{-1}\nabla f(\mathbf{x}_k),
  \]
  siempre que la matriz Hessiana sea invertible \citetext{\citealp{mitnewton2017}; \citealp[Capítulo~7]{nocedal2006}}.

  Si tomamos como punto inicial $\mathbf{x}_0=(0,1)^T$, entonces
  \[
  \nabla f(0,1)=
  \begin{pmatrix}
  0\\
  0
  \end{pmatrix},
  \qquad
  H_f(0,1)=
  \begin{pmatrix}
  0 & \frac{111}{4}\\[2pt]
  \frac{111}{4} & 0
  \end{pmatrix}.
  \]

  Aunque la Hessiana es invertible, no es positiva definida, pues
  sus valores propios son $\lambda_1=\frac{111}{4}$ y
  $\lambda_2=-\frac{111}{4}$. Esto indica que el punto estacionario
  no es un mínimo local, sino un punto de silla.

  Al sustituir en la iteración de Newton obtenemos
  \[
  \mathbf{x}_1
  =\mathbf{x}_0-H_f(\mathbf{x}_0)^{-1}
  \begin{pmatrix}
  0\\
  0
  \end{pmatrix}
  =\mathbf{x}_0.
  \]
  Por tanto, las iteraciones posteriores permanecen en el mismo
  punto:
  \[
  \mathbf{x}_k=(0,1)^T,\qquad k\geq 0.
  \]

  En consecuencia, el método de Newton se estanca en un punto
  estacionario que no es óptimo.
\end{solucion}

  
  \item Por otro lado, si iniciamos en el punto $(0,1)$, ¿qué procedimiento se puede llevar a cabo para que el método Newton o de quasi-Newton continúe con sus búsquedas hasta alcanzar el óptimo?

    \begin{solucion}
    En el punto exacto $(0,1)$, una búsqueda lineal no genera movimiento porque el gradiente es cero. Primero se puede salir del punto de silla perturbando la iteración en la dirección de curvatura negativa, que para esta Hessiana es $\vec{v}=(1,-1)$. En efecto, para $\varepsilon$ suficientemente pequeño,
    \[
      f(\varepsilon,1-\varepsilon)
      =f(0,1)-\frac{111}{4}\varepsilon^2+O(\varepsilon^3)<f(0,1).
    \]
    Desde ese punto se puede continuar con Newton modificado o quasi-Newton, usando una búsqueda lineal con retroceso y modificando la Hessiana o su aproximación para obtener una dirección de descenso. Así se evita aceptar pasos que aumenten la función y se puede seguir iterando hasta satisfacer el criterio de convergencia \citep[Capítulos~4, 7 y 9]{nocedal2006}.
  \end{solucion}

  
  \item Mencione algunas de las principales desventajas y limitantes del método de Newton.

  \begin{solucion}
    Entre sus principales limitantes se encuentran las siguientes \citep[Capítulo~7]{nocedal2006}:
    \begin{itemize}[leftmargin=1.5em]
      \item Requiere calcular la Hessiana y resolver un sistema lineal en cada iteración, lo que puede ser costoso en problemas de muchas variables.
      \item La Hessiana puede ser singular o indefinida, de modo que el paso de Newton no exista o no sea una dirección de descenso.
      \item La convergencia es local y depende del punto inicial; si este se encuentra lejos del óptimo, el método puede divergir o converger a un punto estacionario que no sea mínimo.
    \end{itemize}
  \end{solucion}
\end{enumerate}
    
\end{ejercicio}



\begin{ejercicio}[20 puntos] Sólo en este ejercicio se admite código, aún así es necesario desarrollar analíticamente todo el procedimiento. El costo del petróleo refinado enviado a Japón a través del Estrecho de Malaca en dólares por kilolitro se representa como la suma del costo del petróleo crudo, el seguro, las aduanas, el costo del flete del petróleo, el costo de carga y descarga, el costo del muelle, el costo de la tubería submarina, el costo de almacenamiento, el costo de la zona de tanques, el costo de refinado y el costo de flete de los productos. Más precisamente, la función a minimizar está dada por
\[
\begin{aligned}
  c(q,t)={}&c_c+c_i+c_x
  +\frac{2.09\times10^4\,t^{-0.3017}}{360}
  +\frac{1.064\times10^6\,t^{0.4925}}{52.47\,q\,(360)}\\
  &+\frac{4.242\times10^4\,a\,t^{0.7952}+1.813\,i\,p\,(nt+1.2q)^{0.861}}{52.47\,q\,(360)}\\
  &+\frac{4.25\times10^3\,a\,(nt+1.2q)}{52.47\,q\,(360)}
  +\frac{5.042\times10^3\,q^{-0.1899}}{360}
  +\frac{0.1049\,q^{0.671}}{360},
\end{aligned}
\]
donde
\[
\begin{array}{r @{\;=\;} l @{\qquad\qquad\qquad\quad} l}
  a   & 0.20 & \text{gastos fijos anuales,}\\
  c_c & 12.5 & \text{precio del crudo, \$/kL,}\\
  c_i & 0.50 & \text{costo del seguro, \$/kL,}\\
  c_x & 0.90 & \text{costo de aduana, \$/kL,}\\
  i   & 0.10 & \text{tipo de interés,}\\
  n   & 2    & \text{número de puertos,}\\
  p   & 7000 & \text{precio del terreno, \$/m}^2\text{,}\\
  q   & \multicolumn{2}{l}{\text{capacidad de la refinería, bbl/día,}}\\
  t   & \multicolumn{2}{l}{\text{tamaño del petrolero, kL.}}
\end{array}
\]
Calcule el costo mínimo del petróleo y el tamaño óptimo del petrolero $t$ y de la refinería $q$ usando los métodos de Newton y el método quasi-Newton. Utilice la conversión 1 kL = 6,29 bbl.


\begin{solucion}

El codigo completo de la solución explicado de manera detallada se encuentra en: \href{https://eduardo-alanis-garcia.github.io/MCDI_1er_Semestre/Matematicas/Unidad%203.%20C%C3%A1lculo%20Multivariable/code/Ejercicio%206%20explicado.html}{Ejercicio 6, Tarea 3. Eduardo Alanís García}.



  Sustituimos los valores constantes y escribimos la función en una forma más compacta. Sea $S=nt+1.2q$ y
  \[
  \begin{gathered}
    c_0=c_c+c_i+c_x=13.9,\qquad r=-0.3017,\quad s=0.4925,\quad u=0.7952,\\
    v=0.861,\quad w=-0.1899,\quad m=0.671,
  \end{gathered}
  \]
  junto con
  \[
  \begin{aligned}
    A&=\frac{2.09\times10^4}{360},&
    B&=\frac{1.064\times10^6}{52.47(360)},&
    C&=\frac{4.242\times10^4a}{52.47(360)},\\
    D&=\frac{1.813ip}{52.47(360)},&
    E&=\frac{4.25\times10^3a}{52.47(360)},&
    F&=\frac{5.042\times10^3}{360},\qquad
    G=\frac{0.1049}{360}.
  \end{aligned}
  \]
  Entonces
  \[
    c(t,q)=c_0+At^r+B\frac{t^s}{q}+C\frac{t^u}{q}
      +D\frac{S^v}{q}+E\frac{S}{q}+Fq^w+Gq^m.
  \]
  Se buscan los puntos estacionarios resolviendo $c_t=0$ y $c_q=0$. Aplicando las reglas de potencias y cociente,
  \[
  \begin{aligned}
    c_t={}&Ar t^{r-1}+Bs\frac{t^{s-1}}{q}
      +Cu\frac{t^{u-1}}{q}+Dvn\frac{S^{v-1}}{q}+E\frac{n}{q}=0,\\
    c_q={}&-B\frac{t^s}{q^2}-C\frac{t^u}{q^2}
      +D\left(-\frac{S^v}{q^2}+1.2v\frac{S^{v-1}}{q}\right)\\
      &+E\left(-\frac{S}{q^2}+\frac{1.2}{q}\right)
      +Fwq^{w-1}+Gmq^{m-1}=0.
  \end{aligned}
  \]
  Para aplicar Newton se requiere la Hessiana. Sus entradas, obtenidas derivando nuevamente, son
  \[
  \begin{aligned}
    c_{tt}={}&Ar(r-1)t^{r-2}+Bs(s-1)\frac{t^{s-2}}{q}
      +Cu(u-1)\frac{t^{u-2}}{q}
      +Dv(v-1)n^2\frac{S^{v-2}}{q},\\
    c_{tq}={}&-Bs\frac{t^{s-1}}{q^2}-Cu\frac{t^{u-1}}{q^2}
      +Dvn\left(-\frac{S^{v-1}}{q^2}+1.2(v-1)\frac{S^{v-2}}{q}\right)
      -E\frac{n}{q^2},\\
    c_{qq}={}&2B\frac{t^s}{q^3}+2C\frac{t^u}{q^3}\\
      &+D\left(2\frac{S^v}{q^3}-2.4v\frac{S^{v-1}}{q^2}
        +1.44v(v-1)\frac{S^{v-2}}{q}\right)\\
      &+E\left(2\frac{S}{q^3}-\frac{2.4}{q^2}\right)
      +Fw(w-1)q^{w-2}+Gm(m-1)q^{m-2}.
  \end{aligned}
  \]
  El paso de Newton es
  \[
    \begin{pmatrix}t_{k+1}\\q_{k+1}\end{pmatrix}
    =
    \begin{pmatrix}t_k\\q_k\end{pmatrix}
    -H_c(t_k,q_k)^{-1}\nabla c(t_k,q_k),
    \qquad
    H_c=\begin{pmatrix}c_{tt}&c_{tq}\\c_{tq}&c_{qq}\end{pmatrix}.
  \] \citep[Capítulo~7]{nocedal2006}.
  Para el método quasi-Newton BFGS se actualiza la aproximación inversa de la Hessiana mediante
  \[
    H_{k+1}^{-1}=(I-\rho_k s_k y_k^T)H_k^{-1}(I-\rho_k y_k s_k^T)
      +\rho_k s_k s_k^T,
    \quad
    s_k=z_{k+1}-z_k,\quad y_k=\nabla c(z_{k+1})-\nabla c(z_k),
    \quad \rho_k=\frac{1}{y_k^Ts_k}.
  \] \citep[Capítulo~9]{nocedal2006}.
  En ambos métodos se utilizó una búsqueda lineal de Armijo para conservar $t,q>0$ y disminuir el costo en cada paso \citep[Capítulo~4]{nocedal2006}. Para el cálculo numérico definimos las coordenadas escaladas
  \[
    \tau=\frac{t}{10^5\,\mathrm{kL}},\qquad
    \kappa=\frac{q}{10^5\,\text{bbl/día}}.
  \]
  Así, $t=10^5\tau\,\mathrm{kL}$ y $q=10^5\kappa\,\text{bbl/día}$. Esta elección se debe a que los tamaños relevantes de $t$ y $q$ son del orden de $10^5$ en sus unidades respectivas; al escalarlos, las coordenadas toman valores moderados: el punto inicial $(\tau_0,\kappa_0)=(0.5,1)$ corresponde a $(t_0,q_0)=(50\,000\text{ kL},100\,000\text{ bbl/día})$, y el óptimo queda cerca de $(4.22,2.05)$. Las iteraciones dieron:
  \[
  \begin{array}{c|r|r|r}
    \text{Método}&t\ (\text{kL})&q\ (\text{bbl/día})&c\ (\text{\$/kL})\\\hline
    \text{Newton}&421\,572.2112&204\,504.6473&18.02506914\\
    \text{BFGS}&421\,572.2065&204\,504.6467&18.02506914
  \end{array}
  \]
  Newton convergió en 9 iteraciones y BFGS en 15; en ambos casos la norma del gradiente escalado fue menor que $10^{-9}$. En la solución, los valores propios de la Hessiana escalada son aproximadamente $0.0151$ y $0.2429$, ambos positivos, lo que confirma que se trata de un mínimo local estricto. Este análisis local no demuestra por sí solo que el mínimo sea global. 
  
  Usando $1\text{ kL}=6.29\text{ bbl}$, el tamaño óptimo del petrolero también puede expresarse como
  \[
    t^*=421\,572.21\text{ kL}\times6.29
      =2\,651\,689.21\text{ bbl}.
  \]
  La capacidad óptima de la refinería, convertida a kilolitros por día, es
  \[
    q^*=\frac{204\,504.65}{6.29}=32\,512.66\text{ kL/día}.
  \]
  Por lo tanto, el costo mínimo calculado con la función tal como está escrita en el enunciado es aproximadamente $18.0251\text{ dólares por kL}$, con petrolero de $421\,572.21\text{ kL}$ y refinería de $204\,504.65\text{ bbl/día}$.

Recordemos que el codigo se encuentra de manera detallada en: \href{https://eduardo-alanis-garcia.github.io/MCDI_1er_Semestre/Matematicas/Unidad%203.%20C%C3%A1lculo%20Multivariable/code/Ejercicio%206%20explicado.html}{Ejercicio 6, Tarea 3. Eduardo Alanís García}. Pero el codigo en general fue el siguiente:

\begin{lstlisting}
options(scipen = 999)

# Parámetros del modelo
a = 0.20
c_c = 12.5
c_i = 0.50
c_x = 0.90
i_rate = 0.10
n_ports = 2
p_land = 7000
conversion_bbl_per_kL = 6.29

# Exponentes
r_exp = -0.3017
s_exp = 0.4925
u_exp = 0.7952
v_exp = 0.861
w_exp = -0.1899
m_exp = 0.671

# Coeficientes de la función compacta 
c0 = c_c + c_i + c_x
A = 2.09e4 / 360
B = 1.064e6 / (52.47 * 360)
C = 4.242e4 * a / (52.47 * 360)
D = 1.813 * i_rate * p_land / (52.47 * 360)
E = 4.25e3 * a / (52.47 * 360)
F = 5.042e3 / 360
G = 0.1049 / 360

# Escalas de las variables
scale_t = 1e5 # kL
scale_q = 1e5 # bbl/día

# Función de costo c(t,q), en dólares por kL.
oil_cost = function(t, q) {
  if (length(t) != 1L || length(q) != 1L ||
      !is.finite(t) || !is.finite(q) || t <= 0 || q <= 0) {
    return(Inf)
  }

  S = n_ports * t + 1.2 * q
  c0 +
    A * t^r_exp +
    B * t^s_exp / q +
    C * t^u_exp / q +
    D * S^v_exp / q +
    E * S / q +
    F * q^w_exp +
    G * q^m_exp
}

# Gradiente analítico respecto de (t,q).
oil_gradient = function(t, q) {
  S = n_ports * t + 1.2 * q

  c_t = A * r_exp * t^(r_exp - 1) +
    B * s_exp * t^(s_exp - 1) / q +
    C * u_exp * t^(u_exp - 1) / q +
    D * v_exp * n_ports * S^(v_exp - 1) / q +
    E * n_ports / q

  c_q = -B * t^s_exp / q^2 -
    C * t^u_exp / q^2 +
    D * (-S^v_exp / q^2 + 1.2 * v_exp * S^(v_exp - 1) / q) +
    E * (-S / q^2 + 1.2 / q) +
    F * w_exp * q^(w_exp - 1) +
    G * m_exp * q^(m_exp - 1)

  c(c_t = c_t, c_q = c_q)
}

# Hessiana analítica respecto de (t,q).
oil_hessian = function(t, q) {
  S = n_ports * t + 1.2 * q

  c_tt = A * r_exp * (r_exp - 1) * t^(r_exp - 2) +
    B * s_exp * (s_exp - 1) * t^(s_exp - 2) / q +
    C * u_exp * (u_exp - 1) * t^(u_exp - 2) / q +
    D * v_exp * (v_exp - 1) * n_ports^2 * S^(v_exp - 2) / q

  c_tq = -B * s_exp * t^(s_exp - 1) / q^2 -
    C * u_exp * t^(u_exp - 1) / q^2 +
    D * v_exp * n_ports *
      (-S^(v_exp - 1) / q^2 +
         1.2 * (v_exp - 1) * S^(v_exp - 2) / q) -
    E * n_ports / q^2

  c_qq = 2 * B * t^s_exp / q^3 +
    2 * C * t^u_exp / q^3 +
    D * (2 * S^v_exp / q^3 -
           2.4 * v_exp * S^(v_exp - 1) / q^2 +
           1.44 * v_exp * (v_exp - 1) * S^(v_exp - 2) / q) +
    E * (2 * S / q^3 - 2.4 / q^2) +
    F * w_exp * (w_exp - 1) * q^(w_exp - 2) +
    G * m_exp * (m_exp - 1) * q^(m_exp - 2)

  matrix(c(c_tt, c_tq, c_tq, c_qq), nrow = 2L, byrow = TRUE,
         dimnames = list(c("t", "q"), c("t", "q")))
}

# Transformaciones entre variables físicas (t,q) y escaladas (tau,kappa).

to_physical = function(z) {
  c(t = scale_t * z[1], q = scale_q * z[2])
}

scaled_cost = function(z) {
  physical = to_physical(z)
  oil_cost(physical["t"], physical["q"])
}

scaled_gradient = function(z) {
  physical = to_physical(z)
  c(scale_t, scale_q) *
    oil_gradient(physical["t"], physical["q"])
}

scaled_hessian = function(z) {
  physical = to_physical(z)
  scale_matrix = diag(c(scale_t, scale_q))
  scale_matrix %*%
    oil_hessian(physical["t"], physical["q"]) %*%
    scale_matrix
}

# Búsqueda lineal de Armijo; descarta puntos fuera del dominio t,q > 0.
armijo_search = function(z, direction, gradient, value,
                          c1 = 1e-4, contraction = 0.5,
                          min_step = 2^-50) {
  slope = sum(gradient * direction)
  if (!is.finite(slope) || slope >= 0) {
    return(NULL)
  }

  step = 1
  while (step >= min_step) {
    candidate = z + step * direction
    if (all(is.finite(candidate)) && all(candidate > 0)) {
      candidate_value = scaled_cost(candidate)
      if (is.finite(candidate_value) &&
          candidate_value <= value + c1 * step * slope) {
        return(list(z = candidate, value = candidate_value, step = step))
      }
    }
    step = contraction * step
  }

  NULL
}

newton_method = function(z0 = c(tau = 0.5, kappa = 1),
                          tolerance = 1e-9, max_iterations = 100L) {
  z = as.numeric(z0)

  for (iteration in 0:max_iterations) {
    value = scaled_cost(z)
    gradient = scaled_gradient(z)
    gradient_norm = sqrt(sum(gradient^2))

    if (gradient_norm < tolerance) {
      return(list(z = z, value = value, iterations = iteration,
                  gradient_norm = gradient_norm, converged = TRUE))
    }
    if (iteration == max_iterations) {
      break
    }

    hessian = scaled_hessian(z)
    newton_step = tryCatch(
      as.vector(solve(hessian, gradient)),
      error = function(e) rep(NA_real_, 2L)
    )
    direction = -newton_step

    # Si la Hessiana no produce descenso, usar -gradiente como respaldo.
    if (any(!is.finite(direction)) || sum(gradient * direction) >= 0) {
      direction = -gradient
    }

    search = armijo_search(z, direction, gradient, value)
    if (is.null(search)) {
      stop("Newton: la búsqueda de Armijo no encontró un paso aceptable.")
    }
    z = search$z
  }

  list(z = z, value = scaled_cost(z), iterations = max_iterations,
       gradient_norm = sqrt(sum(scaled_gradient(z)^2)), converged = FALSE)
}

bfgs_method = function(z0 = c(tau = 0.5, kappa = 1),
                        tolerance = 1e-9, max_iterations = 200L) {
  z = as.numeric(z0)
  inverse_hessian = diag(2L)

  for (iteration in 0:max_iterations) {
    value = scaled_cost(z)
    gradient = scaled_gradient(z)
    gradient_norm = sqrt(sum(gradient^2))

    if (gradient_norm < tolerance) {
      return(list(z = z, value = value, iterations = iteration,
                  gradient_norm = gradient_norm, converged = TRUE))
    }
    if (iteration == max_iterations) {
      break
    }

    direction = -as.vector(inverse_hessian %*% gradient)
    if (any(!is.finite(direction)) || sum(gradient * direction) >= 0) {
      inverse_hessian = diag(2L)
      direction = -gradient
    }

    search = armijo_search(z, direction, gradient, value)
    if (is.null(search)) {
      stop("BFGS: la búsqueda de Armijo no encontró un paso aceptable.")
    }

    z_next = search$z
    gradient_next = scaled_gradient(z_next)
    s = z_next - z
    y = gradient_next - gradient
    ys = sum(y * s)

    # Actualización BFGS de la Hessiana inversa aproximada.
    if (is.finite(ys) && ys > 1e-12 * sqrt(sum(s^2) * sum(y^2))) {
      rho = 1 / ys
      identity = diag(2L)
      left = identity - rho * tcrossprod(s, y)
      right = identity - rho * tcrossprod(y, s)
      inverse_hessian = left %*% inverse_hessian %*% right +
        rho * tcrossprod(s, s)
    } else {
      inverse_hessian = diag(2L)
    }

    z = z_next
  }

  list(z = z, value = scaled_cost(z), iterations = max_iterations,
       gradient_norm = sqrt(sum(scaled_gradient(z)^2)), converged = FALSE)
}

report_result = function(method_name, result) {
  if (!result$converged) {
    warning(method_name, " no alcanzó la tolerancia solicitada.")
  }

  physical = to_physical(result$z)
  eigenvalues = eigen(scaled_hessian(result$z),
                       symmetric = TRUE, only.values = TRUE)$values

  cat("\n", method_name, "\n", sep = "")
  cat(sprintf("Iteraciones: %d\n", result$iterations))
  cat(sprintf("tau = %.10f, kappa = %.10f\n", result$z[1], result$z[2]))
  cat(sprintf("t* = %.4f kL (%.4f bbl)\n",
              physical["t"], physical["t"] * conversion_bbl_per_kL))
  cat(sprintf("q* = %.4f bbl/día (%.4f kL/día)\n",
              physical["q"], physical["q"] / conversion_bbl_per_kL))
  cat(sprintf("Costo = %.10f dólares/kL\n", result$value))
  cat(sprintf("Norma del gradiente escalado = %.3e\n",
              result$gradient_norm))
  cat("Valores propios de la Hessiana escalada:",
      paste(sprintf("%.8f", eigenvalues), collapse = ", "), "\n")
  if (all(eigenvalues > 0)) {
    cat("La Hessiana es positiva definida: mínimo local estricto.\n")
  }
}

# Punto inicial del ejercicio,
z0 = c(tau = 0.5, kappa = 1)
newton_result = newton_method(z0)
bfgs_result = bfgs_method(z0)

report_result("Newton", newton_result)
report_result("BFGS", bfgs_result)
\end{lstlisting}

  
\end{solucion}

\end{ejercicio}


\section*{Conclusiones}
En el Ejercicio 1, $(0,0)$ es el único punto crítico y es un mínimo estricto global, aun
cuando la Hessiana es singular. En el Ejercicio 2, el mínimo global está en $(-2,1)$ con
$f=-3$; la búsqueda univariante converge y la
recta que une $P_1$ y $P_2$ pasa por el óptimo. En el Ejercicio 3, $(1/3,1/3,1/3)$ es el
mínimo restringido global, con $h=1/3$. En el Ejercicio 4, el gradiente descendente con
$\alpha=0.1$ converge a $(32/7,22/7)$. En el Ejercicio 5,
$(0,1)$ es un punto silla donde Newton se estanca, por lo que se requiere perturbar la
iteración o modificar la Hessiana. En el Ejercicio 6, Newton (9 iteraciones) y BFGS (15)
coinciden en un mínimo local estricto de aproximadamente $18.0251$ dólares por kL, con
$t^*\approx421\,572.21$ kL y $q^*\approx204\,504.65$ bbl/día. Newton converge en menos
iteraciones pero exige calcular la Hessiana, mientras que BFGS la aproxima; no se demostró
que el mínimo sea global.

\bibliographystyle{unsrtnat}
\bibliography{referencias}

\end{document}
